% planning/plan-style.tex
% Shared house style for the printed weekly lesson-plan document. \input this
% from the plan file, right after \documentclass.
%
% Do not fork this file per document. If a plan needs a different look, the
% question is whether every plan should have it.
%
% The choices here are deliberate:
%   * LANDSCAPE A4. The day table is four columns of prose; portrait crushes
%     the Main Activities column into a ribbon and the page stops being usable
%     at a glance, which is the only thing this document is for.
%   * ONE DAY PER PAGE, always. A day that spills onto a second page is worse
%     than a day set two points smaller, so the table font is deliberately
%     small and the day page ends with \clearpage rather than \newpage.
%   * NO MATH MODE ANYWHERE. This document never loads mathastext, so a single
%     "$-3$" would render in Computer Modern serif against Lato and look
%     half-typeset. Use \tminus, \ttimes and \textsuperscript instead.
%   * Palette and link colours match the lesson decks and the homework packets.
%
% Provides: \tminus \ttimes \sep \playmark, the \daypage environment with
% \subjectrow, \weekoverview \glancerow \linkline, \announcepage \annsec
% \annsub, the hand-rolled contents macros \tocweek \tocline, and \slidelink
% \planlink \hwlink for the three kinds of live link.

% -- Packages ----------------------------------------------------------------
\usepackage[T1]{fontenc}
\usepackage[utf8]{inputenc}
\usepackage{textcomp}
\usepackage[default]{lato}
\usepackage[a4paper,landscape,top=1.25cm,bottom=1.2cm,left=1.4cm,right=1.4cm,%
            headheight=14pt,headsep=10pt,footskip=20pt]{geometry}
\usepackage{xcolor}
\usepackage{colortbl}
\usepackage{array}
\usepackage{pifont}
\usepackage{fancyhdr}
\usepackage{multicol}
\usepackage[most]{tcolorbox}
\usepackage{enumitem}
\usepackage{needspace}
\usepackage{hyperref}

\setlength{\parindent}{0pt}
\setlength{\parskip}{0pt}

% -- The Learner's Book palette (same hex values as decks and homework) -------
\definecolor{cbteal}{HTML}{0087A8}
\definecolor{cbpurple}{HTML}{5C2483}
\definecolor{cborange}{HTML}{C25E12}
\definecolor{cbgreen}{HTML}{4A8B23}
\definecolor{cbred}{HTML}{C8102E}
\definecolor{cbblue}{HTML}{1A5FA8}
\definecolor{rulegray}{HTML}{9AA5AD}
\definecolor{headgray}{HTML}{F2F2F2}

\hypersetup{
  colorlinks=true,
  linkcolor=cbteal,
  urlcolor=cbblue,
  bookmarksopen=true,
  bookmarksnumbered=false,
  pdfstartview={FitH},
  pdfpagelayout=SinglePage,
}

% -- Text-mode substitutes for the three glyphs that would otherwise need maths
\newcommand{\tminus}{\textminus}
\newcommand{\ttimes}{\texttimes}
\newcommand{\sep}{\,\textperiodcentered\,}
\newcommand{\playmark}{\textcolor{cbteal}{\ding{228}}}

% -- The three kinds of live link ---------------------------------------------
% Slides: the deck the class actually saw, on the deployed site.
% Plan:   the one-page teacher plan for the same unit.
% HW:     the homework shelf. Individual packet PDFs are fingerprinted by Vite
%         at build time, so their URLs change on every deploy -- the shelf page
%         is the only stable address for a packet, and it lists them all.
\newcommand{\siteroot}{https://bowenpra.github.io/classroom/\#}
\newcommand{\slidelink}[2]{\playmark~\href{\siteroot/lesson/#1}{Slides \sep #2}}
\newcommand{\planlink}[2]{\href{\siteroot/plan/#1}{Teacher plan \sep #2}}
\newcommand{\hwlink}[1]{\href{\siteroot/homework}{#1}}

% Slides link for the PARENT UPDATE ONLY. Copying a hyperlinked label out of
% the PDF into a Google Classroom post loses the href -- Classroom only
% auto-links plain URL text -- so here the visible text IS the full address,
% not a friendly label. #1 course/unit (the href target); #2 human label
% (printed above the link, for context); #3 the SAME address as #1, typed out
% by hand with \# and \_ escaped, since a macro argument can't be redisplayed
% as literal text without that escaping.
\newcommand{\slidelinkfull}[3]{%
  \needspace{3\baselineskip}%
  \playmark~Slides \sep #2\par\vspace{2pt}%
  \href{\siteroot/lesson/#1}{#3}}

% -- Column widths for the day table -----------------------------------------
% Ratios carried over from the Word version so the two look like one document.
% They sum to 25.1cm, which is \textwidth less tabcolsep and rules.
\newlength{\colSubj}\setlength{\colSubj}{2.5cm}
\newlength{\colLO}  \setlength{\colLO}{6.0cm}
\newlength{\colAct} \setlength{\colAct}{9.4cm}
\newlength{\colEv}  \setlength{\colEv}{7.2cm}

\newcolumntype{L}[1]{>{\raggedright\arraybackslash}p{#1}}

% -- A line inside a table cell ----------------------------------------------
% Cells are lists of short statements, not paragraphs. \li sets one and leaves
% a gap after it, so a cell is written \li{...}\li{...} with no blank lines.
\newcommand{\li}[1]{#1\par\vspace{3.2pt}}

% -- Running head and foot ----------------------------------------------------
\newcommand{\currentweek}{}
\renewcommand{\currentweek}{}
\pagestyle{fancy}
\fancyhf{}
\fancyhead[L]{\footnotesize\color{black!55}Mr Bowen\sep Year 7 Mathematics \& Science}
\fancyhead[R]{\footnotesize\color{black!55}\currentweek}
\fancyfoot[R]{\footnotesize\color{black!55}\thepage}
\fancyfoot[L]{\footnotesize\color{black!40}\hyperref[toc]{Contents}}
\renewcommand{\headrulewidth}{0pt}
\renewcommand{\headrule}{\vspace{2pt}\hbox to\headwidth{\color{cbteal!45}\leaders\hrule height 1.1pt\hfill}}

% -- Page headings ------------------------------------------------------------
% A big coloured title with a thin rule under it. Used by every page type so
% the document reads as one thing.
\newcommand{\planheading}[2]{%
  \noindent{\LARGE\bfseries\color{#1}#2}\par\vspace{4pt}%
  \hbox to \textwidth{\color{#1!35}\leaders\hrule height 1.6pt\hfill}\par\vspace{10pt}}

% The same thing one step down, for a second block on a page that already has a
% \planheading at the top of it.
\newcommand{\planheadsub}[2]{%
  \noindent{\Large\bfseries\color{#1}#2}\par\vspace{3pt}%
  \hbox to \textwidth{\color{#1!35}\leaders\hrule height 1.3pt\hfill}\par\vspace{8pt}}

% -- The day page -------------------------------------------------------------
% One weekday, one page, always. #1 is the label key used by the contents.
\newenvironment{daypage}[3]{%
  \clearpage
  \phantomsection\label{#1}%
  \pdfbookmark[1]{#2 #3}{#1}%
  \planheading{cbteal}{#2 #3}%
  \footnotesize
  \renewcommand{\arraystretch}{1.25}%
  \arrayrulecolor{black!45}%
  \noindent\begin{tabular}{|L{\colSubj}|L{\colLO}|L{\colAct}|L{\colEv}|}
  \hline
  \rowcolor{headgray}
  \centering\textbf{Subject} & \centering\textbf{Learning Objective / Focus} &
  \centering\textbf{Main Activities} & \centering\arraybackslash\textbf{Evidence of Learning} \\
  \hline
}{%
  \end{tabular}\par
}

% \subjectrow{Homeroom}{8:30 -- 8:45}{objective}{activities}{evidence}
\newcommand{\subjectrow}[5]{%
  \textbf{#1}\par\vspace{1pt}{\color{black!55}#2} & #3 & #4 & #5 \\
  \hline}

% -- The week overview page ---------------------------------------------------
\newenvironment{weekoverview}[3]{%
  \clearpage
  \phantomsection\label{#1}%
  \pdfbookmark[0]{#2 \textperiodcentered\ #3}{#1}%
  \planheading{cbpurple}{#2 \sep #3}%
  \small
  \renewcommand{\arraystretch}{1.3}%
  \arrayrulecolor{black!45}%
  \noindent\begin{tabular}{|L{3.1cm}|L{10.7cm}|L{11.3cm}|}
  \hline
  \rowcolor{headgray}
  \centering\textbf{Day} & \centering\textbf{Science \sep 8:45 -- 9:35} &
  \centering\arraybackslash\textbf{Mathematics \sep 10:45 -- 11:35} \\
  \hline
}{%
  \end{tabular}\par\vspace{14pt}%
}

\newcommand{\glancerow}[4]{%
  \hyperref[#1]{\textbf{#2}} & #3 & #4 \\ \hline}

% A block of links under the at-a-glance table. Two columns, because the list
% is short lines and a single column wastes half a landscape page.
\newenvironment{weeklinks}{%
  {\normalsize\bfseries\color{cborange}\ding{228}~Links for this week}\par\vspace{6pt}%
  \small
  \begin{multicols}{2}
}{%
  \end{multicols}
}
\newcommand{\linkline}[1]{#1\par\vspace{4pt}}

% -- The parent announcement --------------------------------------------------
% Deliberately NOT its own page. It follows the at-a-glance table and the links
% on the week page, because it is week-level material and the overview alone
% leaves two thirds of a landscape page empty. A long update simply flows onto
% a second page; the next \daypage clears it either way.
\newcommand{\annstart}[2]{%
  \par\vspace{15pt}%
  \needspace{7\baselineskip}%
  \phantomsection\label{#1}%
  \pdfbookmark[1]{Parent update -- #2}{#1}%
  \planheadsub{cbgreen}{Parent update \sep #2}}

\newenvironment{annbody}{\small\begin{multicols}{2}}{\end{multicols}}

\newcommand{\annsec}[1]{%
  \par\vspace{9pt}%
  {\bfseries\color{cbgreen}\MakeUppercase{#1}}\par\vspace{2pt}%
  \hbox to \columnwidth{\color{cbgreen!35}\leaders\hrule height 1pt\hfill}\par\vspace{6pt}}

\newcommand{\annsub}[1]{%
  \par\vspace{6pt}{\bfseries\color{black!80}#1}\par\vspace{2pt}}

% A note from Mr Bowen to himself, never part of the announcement text.
\newtcolorbox{editorial}[1][Note]{%
  enhanced, breakable, arc=5pt, boxrule=1pt,
  colback=cborange!6, colframe=cborange, coltitle=cborange,
  colbacktitle=white, fonttitle=\bfseries\small,
  attach boxed title to top left={xshift=11pt, yshift=-7pt},
  boxed title style={arc=3pt, boxrule=1pt, left=6pt, right=6pt, top=1pt, bottom=1pt},
  left=10pt, right=10pt, top=12pt, bottom=8pt, title={#1}, fontupper=\small}

% -- The hand-rolled table of contents ----------------------------------------
% Hand-rolled rather than \tableofcontents because this document's structure is
% weeks and days, not sections, and because \tableofcontents inside multicols
% is fragile. Every entry is a live internal link plus a page number.
\newcommand{\tocweek}[2]{%
  \par\vspace{9pt}%
  {\bfseries\color{cbpurple}\hyperref[#1]{#2}}\par\vspace{3pt}}
% The colour must be set OUTSIDE the leader box. xcolor's \color queues an
% \aftergroup\reset@color, so colouring the box from within lands that reset
% between the leaders and the \hfill and TeX stops with "Leaders not followed
% by proper glue" -- once per contents line.
\newcommand{\tocline}[2]{%
  \noindent\hspace{1.1em}\hyperref[#1]{#2}%
  \nobreak\ {\color{black!30}\leaders\hbox to 0.55em{\hss.\hss}\hfill}\ \nobreak
  \hyperref[#1]{\pageref{#1}}\par\vspace{1.5pt}}
